\chapter{Travas}
\label{CH:LOCK}

A maioria dos kernels, incluindo o xv6, intercala a execução
de múltiplas atividades.
Uma fonte de intercalação é o hardware multiprocessado:
computadores com
várias CPUs executando de forma independente, como é o caso do xv6 com RISC-V.
Essas múltiplas CPUs compartilham a RAM física,
e o xv6 explora esse compartilhamento para manter
estruturas de dados que todas as CPUs leem e escrevem.
Esse compartilhamento levanta a possibilidade de
uma CPU ler uma estrutura de dados enquanto outra
CPU está no meio de uma atualização dessa estrutura, ou até mesmo
múltiplas CPUs atualizando os mesmos dados simultaneamente;
sem um projeto cuidadoso, esse acesso paralelo provavelmente
resultará em erros ou em estruturas de dados corrompidas.
Mesmo em um sistema uniprocessado, o kernel pode alternar a CPU entre
várias threads, causando a intercalação de suas execuções.
Finalmente, um manipulador de interrupção de dispositivo que modifica
os mesmos dados que algum código interrompível pode causar danos
se a interrupção ocorrer no momento exato.
A palavra
\indextext{concorrência}
refere-se a situações em que
múltiplos fluxos de instrução são intercalados,
devido ao paralelismo multiprocessado,
troca de threads ou interrupções.

Kernels são repletos de dados acessados concorrentemente. Por exemplo, duas CPUs
podem chamar simultaneamente {\tt kalloc}, acessando concorrentemente
o início da lista de memória livre. Projetistas de kernels gostam de permitir
bastante concorrência, já que isso pode resultar em aumento de desempenho
através do paralelismo e em maior responsividade. No entanto, como consequência,
os projetistas precisam se assegurar da
correção do sistema mesmo diante dessa concorrência. Existem muitas abordagens
para se chegar a um código correto, algumas mais fáceis de raciocinar do que outras. Estratégias
voltadas para garantir a correção sob concorrência, e as abstrações que as sustentam,
são chamadas de técnicas de \indextext{controle de concorrência}.

O xv6 utiliza várias técnicas de controle de concorrência, dependendo da
situação; muitas outras seriam possíveis. Este capítulo foca em uma técnica amplamente
utilizada: a \indextext{trava} (lock). Uma trava fornece exclusão mútua,
garantindo que apenas uma CPU por vez possa possuir a trava. Se
o programador associar uma trava a cada item de dado compartilhado, e o
código sempre adquirir a trava associada antes de acessar o item, então o
item será acessado por apenas uma CPU por vez. Nessa situação, dizemos que
a trava protege o item de dado. Embora travas sejam um mecanismo de controle
de concorrência fácil de entender, sua desvantagem é que podem limitar
o desempenho, pois forçam a serialização de operações concorrentes.

O restante deste capítulo explica por que o xv6 precisa de travas, como o xv6
as implementa e como as utiliza.

% O seguinte é mais relevante para código sem travas:
%
% Uma observação importante é que, ao analisar um trecho de código no
% xv6, é necessário perguntar-se se algum código concorrente poderia
% alterar o comportamento pretendido ao modificar dados (ou recursos de hardware)
% dos quais esse código depende.
% É preciso ter em mente que o compilador pode transformar uma
% única instrução em C em várias instruções de máquina,
% e que essas instruções podem ser executadas de maneira
% intercalada com instruções executando em outras CPUs.
% Ou seja, não se pode assumir que as linhas de código C
% serão executadas de forma atômica.
% A concorrência torna o raciocínio sobre correção mais difícil.

%%
\section{Condições de Corrida (Races)}
%%

\begin{figure}[t]
\center
\includegraphics[scale=0.8]{fig/smp.pdf}
\caption{Arquitetura SMP simplificada}
\label{fig:smp}
\end{figure}

Como exemplo da necessidade de travas, considere dois processos
com filhos terminados chamando
{\tt wait} em dois
CPUs diferentes. {\tt wait} libera a memória do filho.
Assim, em cada CPU, o kernel chamará {\tt kfree}
para liberar as páginas de memória dos filhos. O alocador de memória do kernel mantém uma lista encadeada:
\lstinline{kalloc()} \lineref{kernel/kalloc.c:/^kalloc/} remove
uma página de memória do início da lista de páginas livres, e \lstinline{kfree()}
\lineref{kernel/kalloc.c:/^kfree/} insere uma página de volta na lista.
Para melhor desempenho, pode-se esperar que os {\tt kfree}s dos dois processos pai
executem em paralelo, sem que um precise esperar pelo outro,
mas isso não seria correto dada a implementação de {\tt kfree} no xv6.

A Figura~\ref{fig:smp} ilustra esse cenário com mais detalhes: a
lista encadeada de páginas livres está na memória compartilhada pelos dois CPUs, que
a manipulam usando instruções de leitura e escrita. (Na
realidade, os processadores têm caches, mas conceitualmente
sistemas multiprocessados se comportam como se houvesse uma memória única e compartilhada.)
Se não houvesse requisições concorrentes,
poderíamos implementar uma operação de inserção (\lstinline{push}) na lista como
a seguir:
\begin{lstlisting}[numbers=left,firstnumber=1]
    struct element {
      int data;
      struct element *next;
    };
    
    struct element *list = 0;
    
    void
    push(int data)
    {
      struct element *l;
   
      l = malloc(sizeof *l);
      l->data = data;
      l->next = list; (*@\label{line:next}@*)
      list = l;  (*@\label{line:list}@*)
   }
\end{lstlisting}

\begin{figure}[t]
\center
\includegraphics[scale=0.5]{fig/race.pdf}
\caption{Exemplo de condição de corrida}
\label{fig:race}
\end{figure}

Essa implementação é correta se executada isoladamente.
Entretanto, o código não é correto se mais de uma
instância for executada simultaneamente.
Se dois CPUs executarem
\lstinline{push}
ao mesmo tempo,
ambos podem executar a linha~\ref{line:next} como mostrado na Fig.~\ref{fig:smp},
antes que qualquer um execute a linha~\ref{line:list}, o que resulta
em um resultado incorreto, como ilustrado pela Figura~\ref{fig:race}.
Haveria então dois
elementos da lista com
\lstinline{next}
apontando para o valor antigo de
\lstinline{list}.
Quando as duas atribuições para
\lstinline{list}
ocorrerem na linha~\ref{line:list},
a segunda sobrescreverá a primeira;
o elemento envolvido na primeira atribuição
será perdido.

A atualização perdida na linha~\ref{line:list} é um exemplo de uma
\indextext{condição de corrida (race)}.
Uma race é uma situação na qual uma posição de memória é acessada
concorrentemente, e pelo menos um dos acessos é uma escrita.
Uma race é frequentemente um sinal de bug, seja por uma atualização perdida
(se os acessos forem escritas) ou por uma leitura de
uma estrutura de dados parcialmente atualizada.
O resultado de uma race depende do
código de máquina gerado pelo compilador,
do tempo de execução dos dois CPUs envolvidos,
e de como suas operações de memória são ordenadas pelo sistema de memória,
o que pode tornar erros causados por races difíceis de reproduzir
e depurar.
Por exemplo, adicionar comandos de impressão ao depurar
\lstinline{push}
pode alterar o tempo de execução o suficiente
para que a race desapareça.

A forma usual de evitar races é usar uma trava (lock).
As travas garantem
\indextext{exclusão mútua},
de modo que apenas um CPU por vez possa executar 
as linhas sensíveis de
\lstinline{push};
isso torna o cenário acima impossível.
A versão corretamente sincronizada do código acima
adiciona apenas algumas linhas (destacadas em amarelo):
\begin{lstlisting}[numbers=left,firstnumber=6]
   struct element *list = 0;
   (*@\hl{struct lock listlock;}@*)
    	
   void
   push(int data)
   {
     struct element *l;
     l = malloc(sizeof *l); (*@\label{line:malloc}@*)
     l->data = data;
   
     (*@\hl{acquire(\&listlock);} @*)
     l->next = list;     (*@\label{line:next1}@*)
     list = l;           (*@\label{line:list1}@*)
     (*@\hl{release(\&listlock)}; @*)
   }
\end{lstlisting}
A sequência de instruções entre
\lstinline{acquire}
e
\lstinline{release}
é frequentemente chamada de
\indextext{seção crítica}.
Diz-se que a trava está protegendo
\lstinline{list}.

Quando dizemos que uma trava protege dados, queremos dizer
que ela protege um conjunto de invariantes
relacionados aos dados.
Invariantes são propriedades de estruturas de dados que
são mantidas ao longo das operações.
Tipicamente, o comportamento correto de uma operação depende
dos invariantes serem verdadeiros quando a operação
se inicia. A operação pode violar os invariantes temporariamente,
mas deve restabelecê-los antes de terminar.
Por exemplo, no caso da lista encadeada, o invariante é que
\lstinline{list}
aponta para o primeiro elemento da lista,
e que o campo
\lstinline{next}
de cada elemento aponta para o próximo.
A implementação de
\lstinline{push}
viola esse invariante temporariamente: na linha~\ref{line:next1},
\lstinline{l}
aponta
para o próximo elemento da lista, mas
\lstinline{list}
ainda não aponta para
\lstinline{l}
(esse invariante é restabelecido na linha~\ref{line:list1}).
A condição de corrida examinada acima
ocorreu porque um segundo CPU executou
código que dependia dos invariantes da lista
enquanto eles estavam (temporariamente) violados.
O uso adequado de uma trava garante que apenas um CPU por vez
possa operar sobre a estrutura de dados na seção crítica, de modo que
nenhum CPU executará operações enquanto os invariantes não forem mantidos.

Você pode pensar em uma trava como
\indextext{serializando}
seções críticas concorrentes para que elas rodem uma de cada vez,
preservando os invariantes (assumindo que as seções críticas
estejam corretas isoladamente).
Também é possível pensar que seções críticas protegidas pela mesma trava são
atômicas entre si,
de modo que cada uma vê apenas o conjunto completo de
alterações feitas por seções críticas anteriores, e nunca vê
atualizações parcialmente concluídas.

Apesar de úteis para correção, as travas limitam
o desempenho. Por exemplo, se dois processos chamarem {\tt kfree}
simultaneamente, as travas irão serializar as duas seções críticas,
de modo que não há
benefício em executá-las em diferentes CPUs. Dizemos que múltiplos
processos \indextext{conflitam} se querem a mesma trava ao mesmo
tempo, ou que a trava sofre \indextext{contenção}. Um grande
desafio no projeto de kernels é evitar contenção de travas
na busca por paralelismo. O xv6
faz pouco nesse sentido, mas kernels sofisticados organizam
estruturas de dados e algoritmos especificamente para evitar contenção. No exemplo da lista,
um kernel pode manter uma lista livre separada por CPU e acessar
a lista de outro CPU apenas se a lista local estiver vazia e for necessário "roubar"
memória. Outros casos podem exigir
projetos mais complexos.

A posição das travas também é importante para o desempenho.
Por exemplo, seria correto mover
\lstinline{acquire}
para antes da linha~\ref{line:malloc}
em
\lstinline{push}.
Mas isso provavelmente reduziria o desempenho, pois as chamadas
para
\lstinline{malloc}
seriam serializadas.
A seção ``Usando travas'' abaixo fornece diretrizes sobre onde inserir
as invocações de
\lstinline{acquire}
e
\lstinline{release}.

%%
\section{Código: Travas}
%%

O xv6 possui dois tipos de travas: travas de espera ocupada (spinlocks) e travas de suspensão (sleep-locks).
Vamos começar com as travas de espera ocupada.
O xv6 representa uma spinlock como uma
\indexcode{struct spinlock}
\lineref{kernel/spinlock.h:/struct.spinlock/}.
O campo importante na estrutura é
\lstinline{locked},
uma palavra que vale zero quando a trava está disponível
e diferente de zero quando está em uso.
Logicamente, o xv6 deveria adquirir uma trava executando um código como:

\begin{lstlisting}[numbers=left,firstnumber=21]
   void
   acquire(struct spinlock *lk) // não funciona!
   {
     for(;;) {
       if(lk->locked == 0) {  (*@\label{line:testlocked}@*)
         lk->locked = 1;      (*@\label{line:assign}@*)
         break;
       }
     }
   }
\end{lstlisting}

Infelizmente, essa implementação não garante exclusão mútua em um multiprocessador.
Pode acontecer de dois CPUs simultaneamente
atingirem a linha~\ref{line:testlocked}, verem que 
\lstinline{lk->locked}
é zero, e então ambos adquirem a trava executando a linha~\ref{line:assign}.
Neste ponto, dois CPUs diferentes possuem a trava,
o que viola a propriedade de exclusão mútua.
O que precisamos é de uma maneira de fazer com que as linhas \ref{line:testlocked} e \ref{line:assign}
sejam executadas como uma etapa
\indextext{atômica}
(ou seja, indivisível).

Como as travas são amplamente utilizadas,
processadores multi-core geralmente fornecem instruções que
implementam uma versão atômica das linhas~\ref{line:testlocked} e \ref{line:assign}.
No RISC-V, essa instrução é
\lstinline{amoswap r, a}.
\lstinline{amoswap}
lê o valor no endereço de memória {\tt a},
escreve o conteúdo do registrador {\tt r} nesse endereço,
e coloca o valor lido em {\tt r}.
Ou seja, ela troca o conteúdo do registrador com o da posição de memória.
Ela realiza essa sequência de forma atômica, utilizando hardware especial
para impedir que qualquer outro CPU utilize o mesmo endereço de memória entre a leitura e a escrita.

A função
\indexcode{acquire}
do xv6
\lineref{kernel/spinlock.c:/^acquire/}
usa a chamada da biblioteca C portátil
\lstinline{__sync_lock_test_and_set},
que essencialmente se resume à instrução
\lstinline{amoswap};
o valor de retorno é o conteúdo antigo (trocado) de
\lstinline{lk->locked}.
A função
\lstinline{acquire}
envolve a troca em um laço, repetindo (esperando) até que consiga
adquirir a trava.
Cada iteração troca o valor um com
\lstinline{lk->locked}
e verifica o valor anterior;
se o valor anterior era zero, então a trava foi adquirida,
e a troca terá definido
\lstinline{lk->locked}
como um.
Se o valor anterior era um, então algum outro CPU
já possui a trava, e o fato de termos trocado atômicamente o valor um
em
\lstinline{lk->locked}
não alterou o valor.

Uma vez que a trava é adquirida,
\lstinline{acquire}
registra, para fins de depuração, qual CPU
adquiriu a trava.
O campo
\lstinline{lk->cpu}
é protegido pela trava
e só deve ser modificado enquanto ela estiver adquirida.

A função
\indexcode{release}
\lineref{kernel/spinlock.c:/^release/}
é o oposto de 
\lstinline{acquire}:
ela limpa o campo
\lstinline{lk->cpu}
e então libera a trava.
Conceitualmente, liberar a trava exige apenas atribuir zero a
\lstinline{lk->locked}.
O padrão da linguagem C permite que compiladores implementem uma atribuição
com múltiplas instruções de escrita,
então uma atribuição em C pode não ser atômica com relação a código concorrente.
Em vez disso,
\lstinline{release}
usa a função da biblioteca C
\lstinline{__sync_lock_release}
que realiza uma atribuição atômica.
Essa função também se resume a uma instrução
\lstinline{amoswap}
no RISC-V.

%% 
\section{Código: Usando travas}
%%

O xv6 utiliza travas em muitos lugares para evitar condições de corrida.
Como descrito acima,
\lstinline{kalloc}
\lineref{kernel/kalloc.c:/^kalloc/}
e
\lstinline{kfree}
\lineref{kernel/kalloc.c:/^kfree/}
são um bom exemplo.
Tente os Exercícios 1 e 2 para ver o que acontece se essas
funções omitirem as travas.
Você provavelmente descobrirá que é difícil provocar comportamentos incorretos,
o que sugere que é difícil testar de forma confiável se um código
está livre de erros de sincronização e condições de corrida.
É bem possível que o xv6 ainda contenha corridas não descobertas.

Uma parte difícil ao usar travas é decidir quantas travas
utilizar e quais dados e invariantes cada trava deve proteger.
Há alguns princípios básicos.
Primeiro, sempre que uma variável puder ser escrita por um CPU
ao mesmo tempo em que outro CPU puder lê-la ou escrevê-la,
uma trava deve ser usada para evitar que essas
operações ocorram simultaneamente.
Segundo, lembre-se de que travas protegem invariantes:
se um invariante envolve múltiplas localizações de memória,
tipicamente todas elas devem ser protegidas
por uma única trava para garantir que o invariante seja mantido.

As regras acima dizem quando as travas são necessárias, mas não dizem nada sobre quando elas
são desnecessárias, e isso é importante para a eficiência: travar em excesso
reduz o paralelismo. Se o paralelismo não for importante, então seria possível
simplesmente ter uma única thread e não se preocupar com travas. Um kernel
simples pode fazer isso em um multiprocessador exigindo que uma única trava
seja adquirida ao entrar no kernel e liberada ao sair do kernel (embora
chamadas de sistema bloqueantes como leituras de pipe ou
\lstinline{wait}
possam representar um problema). Muitos sistemas operacionais para uniprocessadores foram adaptados para
funcionar em multiprocessadores usando essa abordagem, às vezes chamada de
``grande trava de kernel'' (big kernel lock), mas essa abordagem sacrifica o paralelismo: apenas um
CPU pode executar no kernel por vez. Se o kernel realizar qualquer computação pesada,
seria mais eficiente usar um conjunto maior de travas mais específicas (fine-grained),
para que o kernel possa ser executado em múltiplos CPUs simultaneamente.

Como exemplo de travamento grosseiro (coarse-grained), o alocador do xv6 em \lstinline{kalloc.c}
possui uma única lista livre protegida por uma única trava. Se
múltiplos processos em diferentes CPUs tentarem alocar páginas ao mesmo
tempo, cada um terá que esperar sua vez girando em {\tt
  acquire}. Esperar (spin) consome tempo de CPU, já que não é trabalho útil.
Se a contenção pela trava estiver desperdiçando uma fração significativa do tempo de CPU,
talvez o desempenho possa ser melhorado alterando o projeto do alocador
para ter múltiplas listas livres, cada uma com sua própria trava, permitindo
alocação realmente paralela.

Como exemplo de travamento refinado (fine-grained), o xv6 possui uma trava separada para
cada arquivo, de forma que processos que manipulam arquivos diferentes geralmente podem
prosseguir sem esperar pelas travas uns dos outros. O esquema de travamento de arquivos
poderia ser ainda mais refinado se quisermos permitir que processos
escrevam simultaneamente em áreas diferentes do mesmo arquivo.
Em última análise, as decisões sobre granularidade das travas devem ser guiadas por medições de desempenho
além de considerações sobre a complexidade do código.

À medida que os próximos capítulos explicarem cada parte do xv6,
eles mencionarão exemplos do uso de travas no xv6
para lidar com concorrência.
Como prévia,
a Figura~\ref{fig:locktable}
lista todas as travas do xv6.

\begin{figure}[t]
\center
\begin{tabular}{ll}
{\bf Trava} & {\bf Descrição} \\
\midrule
bcache.lock & Protege a alocação de entradas no cache de blocos \\
cons.lock & Serializa o acesso ao hardware do console, evita saída embaralhada \\
ftable.lock & Serializa a alocação de uma \texttt{struct file} na tabela de arquivos \\
itable.lock & Protege a alocação de inodes residentes em memória \\
vdisk\_lock & Serializa o acesso ao hardware de disco e à fila de descritores de DMA \\
kmem.lock & Serializa a alocação de memória \\
log.lock & Serializa operações no log de transações \\
pi->lock do pipe & Serializa operações em cada pipe \\
pid\_lock & Serializa os incrementos de \texttt{next\_pid} \\
p->lock do proc & Serializa mudanças no estado de um processo \\
wait\_lock & Ajuda o \texttt{wait} a evitar perdas de notificações (wakeups) \\
tickslock & Serializa operações no contador de \texttt{ticks} \\
ip->lock do inode & Serializa operações em cada inode e seu conteúdo \\
b->lock do buf & Serializa operações em cada bloco de buffer \\
\end{tabular}
\caption{Travas no xv6}
\label{fig:locktable}
\end{figure}

%% 
\section{Deadlock e ordenação de travas}
%%

Se um caminho de execução no kernel precisar manter várias travas ao mesmo tempo, é
importante que todos os caminhos adquiram essas travas na mesma ordem. Caso
contrário, há risco de \indextext{deadlock} (impasse). Suponha que dois caminhos no
xv6 necessitem das travas A e B, mas o caminho 1 as adquire na ordem A
depois B, enquanto o outro caminho as adquire na ordem B depois A.
Imagine que a thread T1 executa o caminho 1 e adquire a trava A,
e a thread T2 executa o caminho 2 e adquire a trava B.
Em seguida, T1 tentará adquirir a trava B, e T2 tentará adquirir a trava A.
Ambas as aquisições irão bloquear indefinidamente, pois em ambos os casos
a outra thread possui a trava necessária, e não a liberará até que
sua chamada de aquisição retorne.
Para evitar tais deadlocks, todos os caminhos devem adquirir
as travas na mesma ordem. A necessidade de uma ordem global para aquisição
de travas implica que as travas são efetivamente parte da especificação de cada função:
quem chama deve invocar as funções de forma que as travas sejam adquiridas
na ordem previamente acordada.

O xv6 possui muitas cadeias de ordenação de travas de comprimento dois envolvendo
travas por processo
(a trava presente em cada
\lstinline{struct proc}),
devido à forma como a função
\lstinline{sleep}
funciona (ver Capítulo~\ref{CH:SCHED}).
Por exemplo,
\lstinline{consoleintr}
\lineref{kernel/console.c:/^consoleintr/}
é a rotina de interrupção que lida com caracteres digitados.
Quando uma nova linha chega, qualquer processo que esteja esperando por
entrada do console deve ser acordado.
Para isso,
\lstinline{consoleintr}
mantém
\lstinline{cons.lock}
enquanto chama 
\indexcode{wakeup},
que adquire 
a trava do processo que está esperando, a fim de acordá-lo.
Como consequência, a ordem global de travas para evitar deadlocks
inclui a regra de que
\lstinline{cons.lock}
deve ser adquirida antes de qualquer trava de processo.
O código do sistema de arquivos contém as maiores cadeias de travas do xv6.
Por exemplo, criar um arquivo requer manter simultaneamente:
uma trava no diretório, uma trava no inode do novo arquivo,
uma trava no buffer de bloco do disco,
a trava do driver de disco \lstinline{vdisk_lock},
e
a trava do processo chamador \lstinline{p->lock}.
Para evitar deadlocks, o código do sistema de arquivos sempre adquire as travas na ordem
mencionada na frase anterior.

Seguir uma ordem global de travas para evitar deadlocks pode ser surpreendentemente
difícil. Às vezes a ordem de travas entra em conflito com a estrutura lógica do
programa; por exemplo, talvez o módulo de código M1 chame o módulo M2, mas a ordem
de travas requer que uma trava em M2 seja adquirida antes de uma trava em M1.
Em outros casos, as identidades das travas não são conhecidas de antemão, talvez
porque uma trava precise ser mantida para descobrir a identidade
da próxima trava a ser adquirida. Esse tipo de situação surge no
sistema de arquivos durante a busca por componentes sucessivos em um nome de caminho, e
no código de {\tt wait} e {\tt exit}, ao percorrer a tabela de
processos em busca de processos filhos. Por fim, o perigo de deadlocks
é frequentemente uma limitação de quão refinado pode ser um esquema de travas,
já que mais travas frequentemente significa mais oportunidade para deadlocks.
A necessidade de evitar deadlocks é frequentemente um fator central na
implementação de kernels.

\section{Travas reentrantes}

Pode parecer que alguns deadlocks e desafios de ordenação de travas poderiam
ser evitados usando \indextext{travas reentrantes}, também chamadas de
\indextext{travas recursivas}. A ideia é que, se uma trava já está sendo mantida
por um processo, e esse processo tentar adquiri-la novamente,
então o kernel poderia simplesmente permitir isso (já que o processo já
possui a trava), em vez de chamar \texttt{panic}, como o kernel do xv6 faz.

No entanto, travas reentrantes tornam mais difícil raciocinar
sobre concorrência: elas quebram a intuição de que travas
fazem com que seções críticas sejam atômicas em relação a outras seções críticas.
Considere as seguintes funções \lstinline{f}
e \lstinline{g}, e uma função hipotética \lstinline{h}:

\begin{lstlisting}
struct spinlock lock;
int data = 0; // protegido por lock

f() {
  acquire(&lock);
  if(data == 0){
    call_once();
    h();
    data = 1;
  }
  release(&lock);
}

g() {
  acquire(&lock);
  if(data == 0){
    call_once();
    data = 1;
  }
  release(&lock);
}

h() {
    ...
}
\end{lstlisting}

Ao observar esse trecho de código, a intuição é que
\lstinline{call_once} será chamada
apenas uma vez: ou por \lstinline{f}, ou por \lstinline{g}, mas não por
ambas.

Porém, se travas reentrantes forem permitidas, e \lstinline{h} acabar chamando
\lstinline{g}, então \lstinline{call_once}
será chamada \emph{duas vezes}.

Se travas reentrantes não forem permitidas, então a chamada de
\lstinline{h} para \lstinline{g}
resultará em um deadlock, o que também não é ótimo. Mas,
supondo que chamar \lstinline{call_once}
duas vezes seja um erro grave, o deadlock
é preferível. O desenvolvedor do kernel observará o deadlock (o
kernel irá entrar em pânico) e poderá corrigir o código para evitá-lo, enquanto
chamar \lstinline{call_once} duas vezes pode
resultar silenciosamente em um erro difícil de depurar.

Por esse motivo, o xv6 usa travas não reentrantes, que são mais simples de entender.
Contudo, desde que os programadores mantenham as regras de travamento em mente,
ambas as abordagens podem funcionar. Se o xv6 fosse usar
travas reentrantes, seria necessário modificar \lstinline{acquire} para
reconhecer que a trava já está sendo mantida pela thread chamadora. Também
seria necessário adicionar um contador de aquisições aninhadas à \lstinline{struct spinlock},
em estilo semelhante ao \lstinline{push_off}, que será discutido a seguir.

\section{Travas e tratadores de interrupção}
\label{s:lockinter}

Algumas spinlocks do xv6 protegem dados que são usados tanto por
threads quanto por tratadores de interrupção.
Por exemplo, o tratador de interrupção do timer
\lstinline{clockintr}
pode incrementar
\indexcode{ticks} 
\lineref{kernel/trap.c:/^clockintr/}
mais ou menos no mesmo momento em que uma thread do kernel
lê
\lstinline{ticks} 
em
\indexcode{sys_sleep}
\lineref{kernel/sysproc.c:/ticks0.=.ticks/}.
A trava
\indexcode{tickslock}
serializa os dois acessos.

A interação entre spinlocks e interrupções gera um perigo potencial.
Suponha que
\indexcode{sys_sleep}
esteja mantendo
\indexcode{tickslock},
e sua CPU seja interrompida por uma interrupção de timer.
\lstinline{clockintr}
tentará adquirir
\lstinline{tickslock},
verá que ela está mantida e aguardará sua liberação.
Nessa situação,
\lstinline{tickslock}
nunca será liberada: apenas
\lstinline{sys_sleep}
pode liberá-la, mas
\lstinline{sys_sleep}
não continuará sua execução até que
\lstinline{clockintr}
retorne.
Portanto, a CPU entrará em deadlock, e qualquer código
que precise dessa trava também ficará congelado.

Para evitar essa situação, se uma spinlock é usada por um tratador de interrupção,
uma CPU nunca deve manter essa trava com interrupções habilitadas.
O xv6 é ainda mais conservador: quando uma CPU adquire qualquer
trava, o xv6 sempre desabilita as interrupções naquela CPU.
Interrupções ainda podem ocorrer em outras CPUs, então
a \lstinline{acquire} de uma interrupção
pode aguardar uma thread liberar uma spinlock — apenas não na mesma CPU.

O xv6 reabilita as interrupções quando a CPU não mantém nenhuma spinlock; ele precisa
fazer um pequeno controle de estado para lidar com seções críticas aninhadas.
\lstinline{acquire}
chama
\indexcode{push_off}
\lineref{kernel/spinlock.c:/^push_off/}
e
\lstinline{release}
chama
\indexcode{pop_off}
\lineref{kernel/spinlock.c:/^pop_off/}
para rastrear o nível de aninhamento de travas na CPU atual.
Quando esse contador chega a zero,
\lstinline{pop_off}
restaura o estado de habilitação de interrupção que existia
no início da seção crítica mais externa.
As funções
\lstinline{intr_off}
e
\lstinline{intr_on}
executam instruções RISC-V para desabilitar e habilitar
interrupções, respectivamente.

É importante que
\indexcode{acquire}
chame
\lstinline{push_off}
estritamente antes de definir
\lstinline{lk->locked}
\lineref{kernel/spinlock.c:/sync_lock_test_and_set/}.
Se a ordem fosse invertida, haveria
uma pequena janela em que a trava
estaria mantida com interrupções habilitadas, e
uma interrupção ocorrendo nesse momento causaria um deadlock no sistema.
De forma similar, é importante que
\indexcode{release}
chame
\indexcode{pop_off}
apenas depois de
liberar a trava
\lineref{kernel/spinlock.c:/sync_lock_release/}.

%% 
\section{Ordenação de instruções e memória}
%% 

É natural pensar que programas executam na ordem
em que as instruções aparecem no código-fonte.
Esse é um modelo mental razoável para código de uma única thread,
mas é incorreto quando
múltiplas threads interagem por meio de memória compartilhada~\cite{riscv:user,boehm04}.
Um dos motivos é que compiladores podem emitir instruções de leitura e escrita
em ordens diferentes das implicadas pelo código-fonte, e podem até mesmo omiti-las
(por exemplo, ao armazenar dados em registradores).
Outro motivo é que a CPU pode executar instruções fora de ordem para melhorar o desempenho.
Por exemplo, uma CPU pode notar que, em uma sequência serial de instruções,
A e B são independentes entre si.
A CPU pode iniciar a instrução B primeiro, ou porque suas
entradas estão prontas antes das de A, ou para sobrepor
a execução de A e B.

Como exemplo do que pode dar errado,
neste código para
\lstinline{push},
seria um desastre se o compilador ou a CPU movesse a
escrita correspondente à
linha~\ref{line:next2} para um ponto após o
\lstinline{release}
na linha~\ref{line:release}:
\begin{lstlisting}[numbers=left,firstnumber=1]
      l = malloc(sizeof *l);
      l->data = data;
      acquire(&listlock);
      l->next = list;   (*@\label{line:next2}@*)
      list = l;      
      release(&listlock);  (*@\label{line:release}@*)
\end{lstlisting}
Se tal reordenação ocorresse, haveria uma janela durante a qual outra CPU
poderia adquirir a trava e
observar o \lstinline{list} atualizado,
mas ver um \lstinline{list->next} não inicializado.

A boa notícia é que compiladores e CPUs ajudam programadores concorrentes
seguindo um conjunto de regras chamado \indextext{modelo de memória}, e
fornecendo alguns mecanismos para ajudar os programadores a controlar a reordenação.

Para informar ao hardware e ao compilador que não devem reordenar,
o xv6 utiliza
\lstinline{__sync_synchronize()} 
tanto em
\lstinline{acquire} \lineref{kernel/spinlock.c:/^acquire/}
quanto em
\lstinline{release} \lineref{kernel/spinlock.c:/^release/}.
\lstinline{__sync_synchronize()}
é uma \indextext{barreira de memória}:
ela instrui o compilador e a CPU a não reordenar leituras ou escritas através da
barreira.
As barreiras em 
\lstinline{acquire}
e
\lstinline{release}
do xv6 impõem a ordem em quase todos os casos em que isso é importante,
já que o xv6 utiliza travas ao redor de acessos a dados compartilhados.
O Capítulo~\ref{CH:LOCK2} discute algumas exceções.

Claro! Aqui está a tradução mantendo a estrutura em LaTeX:

```latex
%% 
\section{Trancas de sono}
%% 

% este material se refere muito ao sleep, yield, etc., então talvez deva
% ser mais tarde, após sleep/wakeup em sched.tex.
% talvez em lock2.tex

Às vezes, o xv6 precisa manter uma tranca por um longo tempo. Por exemplo, o
sistema de arquivos (Chapter~\ref{CH:FS}) mantém um arquivo trancado enquanto lê
e escreve seu conteúdo no disco, e essas operações de disco podem
demorar dezenas de milissegundos. Manter um spinlock por tanto tempo levaria a
desperdício se outro processo quisesse adquiri-lo, já que o processo que está adquirindo
desperdiçaria CPU por um longo tempo enquanto gira. Outro
desvantagem dos spinlocks é que um processo não pode ceder a CPU enquanto
mantém um spinlock; gostaríamos de fazer isso para que outros processos possam
usar a CPU enquanto o processo com a tranca espera pelo disco.
Ceder enquanto mantém um spinlock é ilegal porque pode
levar a um deadlock se um segundo thread tentar adquirir o spinlock;
já que {\tt acquire} não cede a CPU, a rotação do segundo thread pode
impedir que o primeiro thread execute e libere a tranca.
% é verdade que interrupções de timer podem forçar uma troca do
% segundo de volta para o primeiro thread -- mas não se o segundo
% já mantiver uma ou mais trancas.
% é também verdade que o perigo é mais óbvio em um uniprocessador,
% mas também pode surgir com uma cadeia mais longa de threads em
% um multiprocessador.
Ceder enquanto mantém uma tranca também violaria
o requisito de que as interrupções devem estar desativadas enquanto um spinlock é mantido.
Assim, gostaríamos de um tipo de tranca que cede a CPU enquanto espera para
adquirir, e permite yields (e interrupções) enquanto a tranca é mantida.

O xv6 fornece tais trancas na forma de
\indextext{sleep-locks}.
\lstinline{acquiresleep}
\lineref{kernel/sleeplock.c:/^acquiresleep/}
cede a CPU enquanto espera,
usando técnicas que serão explicadas em
Capítulo~\ref{CH:SCHED}.
Em um nível alto, um sleep-locks tem um
campo \lstinline{locked}
que é protegido por um spinlock, e 
a chamada de \lstinline{acquiresleep}
para
\lstinline{sleep}
cede a CPU de forma atômica e libera o spinlock.
O resultado é que outros threads podem executar enquanto
\lstinline{acquiresleep}
espera.

Como os sleep-locks mantêm as interrupções habilitadas, elas não podem ser
usadas em manipuladores de interrupção.
Como
\lstinline{acquiresleep}
pode ceder a CPU,
as trancas de sono não podem ser usadas dentro de seções críticas de spinlock
(embora spinlocks possam ser usados dentro de seções críticas de trancas de sono).

Spin-locks são mais adequados para seções críticas curtas,
já que esperar por elas desperdiça tempo de CPU;
trancas de sono funcionam bem para operações longas.

%% 
\section{Mundo real}
%% 

Programar com trancas continua sendo um desafio, apesar de anos de pesquisa
sobre primitivas de concorrência e paralelismo.
É frequentemente melhor ocultar trancas dentro de
construtos de nível superior, como filas sincronizadas, embora o xv6 não
faça isso. Se você programar com trancas, é prudente usar uma ferramenta que tente
identificar condições de corrida, porque é fácil perder uma invariância que requer
uma tranca.

A maioria dos sistemas operacionais suporta threads POSIX (Pthreads), que permitem que um
processo de usuário tenha várias threads executando concorrentemente em diferentes
CPUs. O Pthreads tem suporte para trancas de nível de usuário, barreiras, etc.
O Pthreads também permite que um programador especifique opcionalmente que uma tranca
deve ser reentrante.

Suportar Pthreads em nível de usuário requer suporte do sistema
operacional. Por exemplo, deve ser o caso que, se um pthread bloqueia
em uma chamada de sistema, outro pthread do mesmo processo deve ser capaz
de executar naquela CPU. Como outro exemplo, se um pthread altera o
espaço de endereços de seu processo (por exemplo, mapeia ou desmapeia memória), o kernel deve
organizar para que outras CPUs que executam threads do mesmo processo atualizem
suas tabelas de páginas de hardware para refletir a mudança no espaço de endereços.

É possível implementar trancas sem instruções atômicas~\cite{lamport:bakery}, mas isso é caro, e a maioria
dos sistemas operacionais usa instruções atômicas.

As trancas podem ser caras se muitas CPUs tentarem adquirir a mesma tranca
ao mesmo tempo. Se uma CPU tem uma tranca
armazenada em seu cache local, e outra CPU deve adquirir a tranca, então a
instrução atômica para atualizar a linha de cache que contém a tranca deve mover a linha
do cache de uma CPU para o cache da outra CPU, e talvez
invalidar quaisquer outras cópias da linha de cache. Buscar uma linha de cache de
outro cache de CPU pode ser ordens de magnitude mais caro do que
buscar uma linha de um cache local.

Para evitar os custos associados às trancas, muitos sistemas operacionais
usam estruturas de dados e algoritmos livres de trancas~\cite{herlihy:art,mckenney:rcuusage}.
Por exemplo, é possível implementar uma lista encadeada como a que está no
início do capítulo que não requer trancas durante as
buscas na lista, e uma instrução atômica para inserir um item em uma lista.
Programar de forma livre de trancas é, no entanto, mais complicado do que programar
com trancas; por exemplo, deve-se se preocupar com a reordenação de instruções e memória.
Programar com trancas já é difícil, então o xv6 evita a
complexidade adicional da programação livre de trancas.

%% 
\section{Exercícios}
%% 

\begin{enumerate}
  
\item Comente as chamadas para
\lstinline{acquire}
e
\lstinline{release}
em
\lstinline{kalloc}
\lineref{kernel/kalloc.c:/^kalloc/}.
Isso parece que deve causar problemas para
o código do kernel que chama
\lstinline{kalloc};
quais sintomas você espera ver?
Quando você executa o xv6, você vê esses sintomas?
E ao executar
\lstinline{usertests}?
Se você não vê um problema, por que não?
Veja se você consegue provocar um problema inserindo
loops fictícios na seção crítica de
\lstinline{kalloc}.

\item Suponha que você em vez disso tenha comentado a
tranca em
\lstinline{kfree} 
(depois de restaurar a tranca em
\lstinline{kalloc}).
O que pode dar errado agora? A falta de trancas em
\lstinline{kfree}
é menos prejudicial do que em
\lstinline{kalloc}?

\item Se duas CPUs chamarem
\lstinline{kalloc}
ao mesmo tempo, uma terá que esperar pela outra,
o que é ruim para o desempenho.
Modifique 
\lstinline{kalloc.c}
para ter mais paralelismo, de modo que chamadas simultâneas
para
\lstinline{kalloc}
de diferentes CPUs possam prosseguir sem esperar uma pela outra.

\item Escreva um programa paralelo usando threads POSIX, que é suportado na maioria
dos sistemas operacionais. Por exemplo, implemente uma tabela hash paralela e meça se
o número de puts/gets escala com o aumento do número de CPUs.

\item Implemente um subconjunto de Pthreads no xv6. Ou seja, implemente uma biblioteca de
threads em nível de usuário para que um processo de usuário possa ter mais de 1 thread e organize
para que essas threads possam ser executadas em paralelo em diferentes CPUs. Crie um
design que trate corretamente uma thread fazendo uma chamada de sistema bloqueante e
mudando seu espaço de endereços compartilhado.

\end{enumerate}

% LocalWords:  CPUs uniprocessor interruptible allocator kalloc kfree
% LocalWords:  struct malloc sizeof listlock invariants spinlock lk
% LocalWords:  amoswap cpu bcache ftable itable inode vdisk DMA kmem
% LocalWords:  pid proc's tickslock inode's ip buf's proc SCHED sys
% LocalWords:  consoleintr wakeup clockintr intr acquiresleep POSIX
% LocalWords:  Pthreads pthread unmaps TLB IPIs
